\documentclass[sigconf,nonacm]{acmart}

%% -----------------------------------------------------------------------
%% Packages
%% -----------------------------------------------------------------------
\usepackage{booktabs}
\usepackage{graphicx}
\usepackage{amsmath}
\usepackage{amssymb}
\usepackage{xspace}
\usepackage{hyperref}
\usepackage{listings}
\usepackage{xcolor}
\usepackage{microtype}

%% -----------------------------------------------------------------------
%% Macros
%% -----------------------------------------------------------------------
\newcommand{\sysname}{\textsc{TemporalEngine}\xspace}
\newcommand{\wcc}{\textsc{WCC}\xspace}
\newcommand{\eps}{\varepsilon}
\newcommand{\lap}{\mathrm{Lap}}

%% -----------------------------------------------------------------------
%% Metadata
%% -----------------------------------------------------------------------
\title{Adversarial Robustness and Privacy-Utility Tradeoffs\\
       in Real-Time Graph-Based Fraud Ring Detection}

\author{Ovindu Karunaratne}
\affiliation{%
  \institution{Independent Research}
  \country{}
}
\email{karunaratneovindu@gmail.com}

\begin{document}

\begin{abstract}
Weakly Connected Component (\wcc) graph analysis is the dominant technique
for real-time fraud ring detection in streaming financial systems.  Prior
adversarial studies of community detection \cite{li2020adversarial,jia2020certified}
target learned models on static social graphs; the streaming, rule-based,
multi-relational \wcc deployments used in production fraud pipelines remain
largely unstudied.  We present \sysname, a production-grade fraud detection
platform built on Apache Flink, Memgraph, and LangGraph, and use it as an
experimental testbed for two questions:
(1)~How does \wcc-based ring detection degrade under adversarial graph
manipulation when the attacker operates transaction-by-transaction across
three independent connection channels (devices, IPs, merchants)?
(2)~What is the privacy-utility tradeoff when Laplace differential
privacy is applied to the co-located velocity detector?

Our experiments (30 Monte Carlo trials per configuration) show that \wcc
exhibits strong \emph{multi-channel redundancy}: breaking only device/IP
connections does not reduce recall ($R = 1.000$ at all budgets), and
breaking only merchant connections is similarly ineffective ($F_1 \geq
0.989$).  Only \emph{simultaneous} evasion of all three connection channels
(devices, IPs, merchants) achieves meaningful recall degradation---and even
then, an adversary must corrupt 80\% of ring members before recall falls
below $0.1$.  For differential privacy, we identify $\eps \approx 5$ as the
practical operating point: below this threshold false-positive rates become
operationally unacceptable ($\mathrm{FPR} = 0.28$ at $\eps=1$); above it,
near-perfect recall is preserved ($\mathrm{FPR} = 0.027$ at $\eps=5$).
\end{abstract}

\keywords{fraud detection, graph analytics, weakly connected components,
          adversarial robustness, differential privacy, streaming systems}

\maketitle

%% ===================================================================
\section{Introduction}
%% ===================================================================

Fraud rings---coordinated groups of accounts sharing devices, IP addresses,
and merchant relationships---account for a disproportionate fraction of
financial fraud losses~\cite{pourhabibi2020}.  Graph-based detection using
Weakly Connected Components (\wcc) is a well-established approach
\cite{akoglu2015,liu2021,liu2018heterogeneous}, adopted by major payment
processors and implemented in production systems such as AWS Graph Fraud
Detection and Neo4j Fraud Detection~\cite{weber2019}.

Despite widespread deployment, two questions remain underexplored:

\textbf{Q1 (Adversarial Robustness).}  How many adversarial modifications
must an attacker make to evade \wcc-based detection, and which connection
channels are the weakest links?  Adversarial work on learned graph models
\cite{zugner2018,dai2018,kim2024fraudgangs} and on community detection in
social networks~\cite{li2020adversarial,jia2020certified} operates on static
graphs and homogeneous edge types.  We study the \emph{streaming,
multi-relational, rule-based} case characteristic of production fraud
pipelines, where an attacker acts transaction-by-transaction and must
simultaneously evade redundant connection channels.

\textbf{Q2 (Privacy-Utility Tradeoff).}  Velocity detection---flagging
users with anomalous transaction rates or volumes in a sliding window---is
the complementary real-time signal to \wcc.  If a financial institution
wants to publish aggregate velocity statistics under differential privacy
guarantees~\cite{dwork2014}, what $\eps$ is needed before the Laplace noise
degrades operational detection quality?

We answer both questions experimentally using \sysname, a streaming fraud
detection system we build and open-source.  Our contributions are:

\begin{enumerate}
  \item A reproducible simulation harness replicating a production \wcc
        fraud ring detector over a multi-relational graph
        (User--Device--IP--Merchant), runnable without a Docker/Kafka stack.
  \item Five channel-aware adversarial strategies (single-channel,
        multi-channel, and orthogonal velocity-throttle attacks), evaluated
        across the full budget range $B \in [0,N]$ and revealing a
        \emph{multi-channel redundancy} effect absent from prior single-edge-type
        community-detection attacks~\cite{li2020adversarial}.
  \item A differential privacy velocity detector applied to \emph{inference-time
        streaming aggregates}---a setting distinct from the training-time
        DP-SGD focus of prior privacy-preserving fraud work---with empirical
        privacy-utility curves across $\eps \in \{0.1, \ldots, 50\}$.
  \item An empirical defense comparison: \wcc baseline vs.\ AdaptiveThreshold
        vs.\ MultiSignalDetector, demonstrating that the vulnerability is
        \emph{structural} once all redundant channels are severed.
\end{enumerate}

%% ===================================================================
\section{System Model}
%% ===================================================================

\subsection{\sysname Architecture}

\sysname processes financial transactions as a streaming pipeline:

\begin{center}
\small
Kafka/Redpanda $\to$ Apache Flink $\to$
$\begin{cases}
  \text{Redis (feature store)}\\
  \text{Memgraph (graph DB)}\\
  \text{Apache Iceberg (data lake)}
\end{cases}$
$\to$ LangGraph Agent
\end{center}

\noindent\textbf{Flink pipeline stages.}
\texttt{SafeParse} $\to$ \texttt{TemporalJoin} $\to$ \texttt{VelocityDetector}
(5-min sliding window) $\to$ \texttt{ContractEnforcer} $\to$ \texttt{PIIShield}
(Microsoft Presidio) $\to$ \texttt{GraphSyncSink} $\to$ \texttt{IcebergSink}.

\noindent\textbf{Graph model.}
Four node types: \texttt{User}, \texttt{Merchant}, \texttt{Device}, \texttt{IP}.
Two edge types: \texttt{TRANSACTED\_WITH} (User$\to$Merchant) and
\texttt{LOGGED\_IN\_FROM} (User$\to$Device, User$\to$IP).

\noindent\textbf{Ring detector (\texttt{ring\_hunter.py}).}
Runs Memgraph's \wcc algorithm periodically.  Components with more than 3
nodes are flagged as candidate fraud rings and dispatched to a LangGraph
investigator agent for autonomous risk decision (ALLOW / BLOCK / INVESTIGATE).

\noindent\textbf{Velocity detector (\texttt{processor.py}).}
Flags a user as \texttt{HIGH} if, within a 5-minute event-time window, they
exceed $\tau_{\mathrm{count}} = 5$ transactions \emph{or} $\tau_{\mathrm{sum}}
= \$1{,}000$.

\subsection{Graph Construction}

For user $u$ transacting at merchant $m$ via device $d$ and IP address
$\mathit{ip}$, we add nodes $\{u, m, d, \mathit{ip}\}$ and edges
$\{(u,m),\,(u,d),\,(u,\mathit{ip})\}$ to a property graph $G$.  A fraud
ring of $k$ users sharing $s$ devices, $t$ IP addresses, and $p$ merchants
forms a \wcc component of size $k + s + t + p$, with all ring members
transitively reachable through any shared infrastructure node.

%% ===================================================================
\section{Threat Model}
%% ===================================================================

We adopt a \textbf{white-box adversarial model}: the attacker knows the
\wcc detection algorithm, all threshold parameters, and the full ring
structure.  This is the strongest reasonable assumption and serves as an
upper bound on adversarial capability.

\noindent\textbf{Attacker goal.}
Minimize the probability of the fraud ring being detected by the \wcc ring
hunter.

\noindent\textbf{Adversarial budget $B$.}
The number of ring users whose transaction parameters the attacker may
modify before being detected by out-of-band controls.  $B=0$ is the benign
baseline; $B=N$ (where $N$ is total ring membership) is the maximum
capability.

\noindent\textbf{Attack surface.}
Three independent connection channels through which users are co-located in
the graph: (1) shared device IDs, (2) shared IP addresses, (3) shared
merchant relationships.

%% ===================================================================
\section{Attack Strategies}
%% ===================================================================

\subsection{DeviceIP Rotation}
Replace targeted users' shared device and IP with unique values.  Merchant
connections remain intact.  \emph{Hypothesis:} merchant co-occurrence alone
maintains ring connectivity, so recall will not decrease.

\subsection{Merchant Rotation}
Replace targeted users' merchant with a unique one-off entity.  Device/IP
sharing remains intact.  \emph{Hypothesis:} device/IP paths alone are
sufficient; merchant rotation will have minimal impact.

\subsection{Full Evasion}
Simultaneously rotate device, IP, \emph{and} merchant for each targeted
user---the maximum adversarial capability.  \emph{Hypothesis:} this severs
all \wcc paths for targeted users; recall approaches zero as $B \to N$.

\subsection{Velocity Throttle (Complementary Attack)}
Splits each targeted user's transaction amount to remain below
$\tau_{\mathrm{sum}}$ per window, defeating the Flink velocity detector
without modifying graph topology.  \emph{Hypothesis:} no effect on \wcc
ring detection; orthogonal to the graph signal.

%% ===================================================================
\section{Defenses}
%% ===================================================================

\subsection{Baseline: Fixed-Threshold WCC}
Component size $> 3$ (\texttt{ring\_hunter.py} default).  Evaluated as
reference across all experiments.

\subsection{AdaptiveThreshold}
Dynamically lowers the component-size cutoff from 3 toward a minimum floor
of 1 as system pressure (recent alert rate) increases.  Increases sensitivity
under suspected attack conditions at the cost of more false positives.

\subsection{MultiSignalDetector}
Requires a component to satisfy at least 2 of 3 independent signals:
\begin{itemize}
  \item Signal A: component size $> 3$
  \item Signal B: device/IP reuse ratio $> 0.3$
        (fraction of infrastructure nodes shared by ${\geq}2$ users)
  \item Signal C: user count in component $> 3$
\end{itemize}
An attacker who breaks Signal~A by rotating infrastructure still triggers
Signal~B or C if sufficient sharing remains.

%% ===================================================================
\section{Differential Privacy for Velocity Detection}
%% ===================================================================

We protect velocity statistics using the Laplace
mechanism~\cite{dwork2014}.  Let $f(D)$ denote a statistic computed over
dataset $D$.  The mechanism outputs $\tilde{f}(D) = f(D) + \eta$ where
$\eta \sim \lap(\Delta f / \eps)$ and $\Delta f$ is the global $L_1$
sensitivity.

\noindent\textbf{Sensitivity analysis.}
\begin{align}
  \Delta_{\mathrm{count}} &= 1
    \quad\text{(adding/removing one tx changes count by $\leq 1$)} \\
  \Delta_{\mathrm{sum}} &= \$500
    \quad\text{(maximum single-transaction amount)}
\end{align}

Noise $\eta_{\mathrm{count}} \sim \lap(1/\eps)$ and
$\eta_{\mathrm{sum}} \sim \lap(500/\eps)$ are independently sampled and
added to \texttt{tx\_count\_5m} and \texttt{tx\_sum\_5m} before threshold
comparison.

We generate $N = 400$ synthetic user windows with known ground-truth labels
(200 \texttt{HIGH}, 200 \texttt{NORMAL}) and report TPR, FPR, FNR across
$\eps \in \{0.1, 0.5, 1.0, 2.0, 5.0, 10.0, 50.0\}$ averaged over 30
independent noise draws.

%% ===================================================================
\section{Experimental Setup}
%% ===================================================================

All experiments use Monte Carlo simulation ($n = 30$ independent trials per
data point, independently seeded).  Each trial instantiates a fresh
\texttt{Scenario} with the parameters in Table~\ref{tab:setup}, applies the
chosen attack, rebuilds the graph from scratch, and evaluates \wcc detection
against ground-truth rings using Intersection-over-Union matching at
$\mathrm{IoU} \geq 0.5$.

\begin{table}[h]
\caption{Experiment parameters}
\label{tab:setup}
\centering
\begin{tabular}{lc}
\toprule
Parameter & Value \\
\midrule
Fraud rings & 10 \\
Users per ring (total: 50) & 5 \\
Shared devices per ring & 2 \\
Shared IPs per ring & 2 \\
Shared merchants per ring & 2 \\
Background (benign) users & 30 \\
IoU threshold (ring match) & 0.5 \\
Budget range ($B$) & 0--50 users (0\%--100\%) \\
Monte Carlo trials & 30 \\
\bottomrule
\end{tabular}
\end{table}

%% ===================================================================
\section{Results}
%% ===================================================================

\subsection{Adversarial Robustness}

Table~\ref{tab:adversarial} presents \wcc detection metrics under each
attack strategy.  Budget is expressed as the percentage of total fraud users
whose transaction parameters are modified.

\begin{table}[h]
\caption{WCC detection metrics under adversarial attack ($n=30$ trials).}
\label{tab:adversarial}
\centering
\small
\begin{tabular}{llccc}
\toprule
Strategy & Budget & Precision & Recall & $F_1$ \\
\midrule
Baseline (no attack) & 0\%   & 0.997 & 1.000 & 0.998 \\
\midrule
DeviceIP Rotation    & 10\%  & 0.985 & 1.000 & 0.992 \\
DeviceIP Rotation    & 40\%  & 0.647 & 1.000 & 0.784 \\
DeviceIP Rotation    & 100\% & 0.615 & 1.000 & 0.760 \\
\midrule
Merchant Rotation    & 10\%  & 0.997 & 1.000 & 0.998 \\
Merchant Rotation    & 100\% & 0.979 & 1.000 & 0.989 \\
\midrule
Full Evasion         & 10\%  & 0.993 & 0.993 & 0.993 \\
Full Evasion         & 40\%  & 0.731 & 0.660 & 0.692 \\
Full Evasion         & 60\%  & 0.516 & 0.313 & 0.388 \\
Full Evasion         & 80\%  & 0.247 & 0.050 & 0.082 \\
Full Evasion         & 100\% & 0.000 & 0.000 & 0.000 \\
\midrule
Velocity Throttle    & 100\% & 0.997 & 1.000 & 0.998 \\
Combined Evasion     & 80\%  & 0.247 & 0.050 & 0.082 \\
\bottomrule
\end{tabular}
\end{table}

\noindent\textbf{F1: Merchant co-occurrence is the dominant \wcc signal.}
DeviceIP rotation degrades precision but never reduces recall below
$1.000$, even when 100\% of ring users rotate devices and IPs.  Fraud rings
remain fully connected through shared merchant relationships, revealing
merchant co-occurrence as the primary structural signal.

\noindent\textbf{F2: Device/IP sharing is redundantly sufficient.}
Merchant rotation is similarly ineffective ($F_1 \geq 0.989$ at all
budgets): device and IP sharing alone sustain full connectivity.

\noindent\textbf{F3: Full Evasion produces a critical transition between
60--80\% budget.}  $F_1$ falls from 0.998 at 0\% budget to 0.692 at 40\%,
0.082 at 80\%, and 0.000 at 100\%.  The inflection point near 60--80\%
budget ($R$ drops from 0.313 to 0.050) indicates that a ring can tolerate
approximately 3 of 5 members being fully evaded before its \wcc footprint
collapses.

\noindent\textbf{F4: Velocity throttle is orthogonal to WCC.}
VelocityThrottle ($100\%$ budget) leaves detection unchanged ($F_1 = 0.998$),
confirming the two detection layers are independent.  Combined Evasion
matches Full Evasion alone.

\noindent\textbf{F5: Multi-channel redundancy is the core defense.}
Partial evasion through any single channel fails to reduce recall;
simultaneous evasion of all three channels is required.

\subsection{Differential Privacy}

Table~\ref{tab:dp} shows velocity detection metrics under Laplace noise.

\begin{table}[h]
\caption{Velocity detection under Laplace DP noise ($n=30$ trials,
         $N=400$ windows).}
\label{tab:dp}
\centering
\begin{tabular}{ccccc}
\toprule
$\eps$ & TPR & FNR & FPR & Accuracy \\
\midrule
0.1  & 0.8598 & 0.1402 & 0.6732 & 0.5933 \\
0.5  & 0.9692 & 0.0308 & 0.4265 & 0.7713 \\
1.0  & 0.9907 & 0.0093 & 0.2782 & 0.8562 \\
2.0  & 0.9988 & 0.0012 & 0.1200 & 0.9394 \\
5.0  & 1.0000 & 0.0000 & 0.0270 & 0.9865 \\
10.0 & 1.0000 & 0.0000 & 0.0097 & 0.9952 \\
50.0 & 1.0000 & 0.0000 & 0.0002 & 0.9999 \\
\bottomrule
\end{tabular}
\end{table}

\noindent\textbf{F6: Phase transition at $\eps \approx 2$--$5$.}
Below $\eps = 2$, FPR exceeds 12\%, rendering the detector operationally
unreliable at production scale.  Above $\eps = 5$, TPR reaches $1.000$ and
FPR drops below 3\%.

\noindent\textbf{F7: 10$\times$ FPR improvement for $5\times$ privacy cost.}
FPR drops from 27.8\% ($\eps=1$) to 2.7\% ($\eps=5$)---a $10\times$
improvement for only a $5\times$ increase in privacy budget.

\noindent\textbf{F8: Recommended operating points.}
$\eps = 2$ is the minimum viable threshold for automated blocking systems
where follow-up investigation is cheap.  $\eps = 5$ is recommended for
human-in-the-loop investigation queues.

\subsection{Defense Comparison}

Table~\ref{tab:defense} compares the three detectors under Full Evasion.

\begin{table}[h]
\caption{Detector $F_1$ under Full Evasion attack ($n=30$ trials).}
\label{tab:defense}
\centering
\begin{tabular}{lcccccc}
\toprule
Detector & 0\% & 20\% & 40\% & 60\% & 80\% & 100\% \\
\midrule
WCC Baseline      & 0.998 & 0.995 & 0.692 & 0.388 & 0.082 & 0.000 \\
AdaptiveThreshold & 0.998 & 0.995 & 0.692 & 0.388 & 0.082 & 0.000 \\
MultiSignal       & 0.998 & 0.995 & 0.700 & 0.397 & 0.083 & 0.000 \\
\bottomrule
\end{tabular}
\end{table}

\noindent\textbf{F9: No defense meaningfully improves WCC robustness.}
MultiSignal offers marginal gains ($F_1 + 0.8\%$ at 40\% budget) but does
not alter the recall degradation curve.  At 100\% budget, all three
detectors fail completely.  This implies the vulnerability is structural:
once all three connection channels are severed, no graph-based single-snapshot
detector can recover the ring.

\noindent\textbf{Implication.}
Effective defense requires either (a) \emph{temporal graph analysis}---detecting
coordinated changes in connection patterns over time rather than evaluating a
single \wcc snapshot---or (b) \emph{non-graph signals} such as behavioral
biometrics or device fingerprinting at the application layer.

%% ===================================================================
\section{Discussion}
%% ===================================================================

\subsection{Cost of Evasion}

For a practical 5-user ring, an adversary must rotate device IDs, IP
addresses, \emph{and} merchant relationships for all members before \wcc
detection fails.  In practice: device ID rotation requires new physical
hardware or emulators (\textasciitilde\$50--200/device); IP rotation is
achievable via VPN or residential proxies (\textasciitilde\$5--20/month/IP);
merchant rotation requires coordination with additional shell merchants---the
highest operational cost.  Financial institutions should accordingly
prioritize the integrity of merchant-side graph edges.

\subsection{Temporal Window Constraints as a Defense}

Our results motivate extending \wcc with \emph{temporal window constraints}:
running \wcc only on edges created within a rolling 7- or 30-day window
would force adversaries to maintain evasion continuously rather than
executing a one-time rotation.  The simulation harness we release makes this
straightforward to evaluate; we leave it as future work.

\subsection{Limits of Laplace DP for Velocity Detection}

The large sensitivity of $\texttt{tx\_sum}$ ($\Delta = \$500$) relative to
the detection threshold (\$1{,}000) explains why DP noise dominates at
low~$\eps$.  A more efficient approach would apply DP to normalized
statistics (e.g., $\texttt{tx\_sum} / \Delta_{\mathrm{sum}} \in [0,1]$,
sensitivity = $1/N$ for $N$ contributing users) or use smooth sensitivity
\cite{nissim2007}.  Either approach would allow significantly tighter privacy
budgets at the same utility level.

%% ===================================================================
\section{Related Work}
%% ===================================================================

\noindent\textbf{Graph-based fraud detection.}
Akoglu et al.~\cite{akoglu2015} survey graph anomaly detection including
\wcc-based methods.  Pourhabibi et al.~\cite{pourhabibi2020} provide a
systematic review of graph-based fraud detection approaches.  Weber et
al.~\cite{weber2019} apply GCNs to financial forensics on the Elliptic
Bitcoin dataset; Liu et al.~\cite{liu2018heterogeneous} deploy heterogeneous
graph neural networks for malicious account detection at Alibaba scale; Liu
et al.~\cite{liu2021} address class imbalance in GNN fraud detection.  \wcc
and community detection remain the most widely deployed production methods
due to their simplicity and interpretability.

\noindent\textbf{Adversarial attacks on graphs.}
Z\"{u}gner et al.~\cite{zugner2018} introduce targeted adversarial
perturbations for GNNs on node classification tasks.  Dai et
al.~\cite{dai2018} formulate graph attacks as a reinforcement learning
problem.  Kim et al.~\cite{kim2024fraudgangs} extend these ideas to
GNN-based fraud detection with multi-target injection attacks.  All three
operate on \emph{learned classifiers} and assume the attacker can observe
or perturb a static graph snapshot.

\noindent\textbf{Adversarial attacks on community detection.}
Closest to our threat model, Li et al.~\cite{li2020adversarial} introduce
adversarial attacks against community detection by hiding individuals in
social graphs.  Jia et al.~\cite{jia2020certified} provide certified
robustness bounds via randomized smoothing.  Our work differs in three
key respects: (i)~we target \wcc applied to a \emph{multi-relational}
graph (Users, Devices, IPs, Merchants) rather than a single edge type,
exposing a redundancy effect absent from single-relation analyses;
(ii)~we operate in a \emph{streaming, production} setting where each
adversarial action is a real transaction subject to the live Flink
pipeline; and (iii)~we evaluate alongside a co-located velocity detector,
quantifying the orthogonality between graph-based and aggregate-based
detection layers.

\noindent\textbf{Differential privacy in ML systems.}
Dwork and Roth~\cite{dwork2014} establish the algorithmic foundations of DP.
Nissim et al.~\cite{nissim2007} introduce smooth sensitivity for tighter
noise calibration.  Prior work applying DP to fraud classifiers focuses on
training-time privacy; we apply DP to \emph{inference-time streaming
aggregates}, a distinct and more practically relevant deployment scenario.

%% ===================================================================
\section{Conclusion}
%% ===================================================================

We present the first adversarial robustness study of \wcc-based fraud ring
detection in a \emph{streaming, multi-relational, production} setting.
Where prior community-detection adversarial work~\cite{li2020adversarial,
jia2020certified} targets learned models on single-edge-type static graphs,
our experiments expose a structural property unique to deployed \wcc
pipelines: \emph{multi-channel redundancy} across heterogeneous edge types.

\begin{enumerate}
  \item \textbf{\wcc is surprisingly robust} due to multi-channel redundancy.
        Breaking any single connection channel (device/IP \emph{or} merchant)
        is insufficient to reduce recall below 1.000.
  \item \textbf{Full evasion requires high adversarial cost.}  An attacker
        must break all three connection channels for all ring members;
        merchant evasion is particularly costly operationally.
  \item \textbf{Differential privacy is deployable at $\eps \approx 5$}
        for velocity detection: $\mathrm{TPR} = 1.0$, $\mathrm{FPR} = 0.027$.
  \item \textbf{No existing graph defense} meaningfully improves \wcc
        robustness under full evasion.  Temporal graph analysis is the most
        promising direction for future work.
\end{enumerate}

All code, simulation harness, and results CSVs are available at
\url{https://github.com/Cookie-Cat21/Temporal-Feature-Engine}.

%% ===================================================================
\bibliographystyle{ACM-Reference-Format}
\bibliography{refs}

\end{document}
